%%%%% This template was made by inexcel (email: inexcel98@163.com) for the presentation at his master's defense.
%%%%% The related elements are from Southeast University (https://www.seu.edu.cn/bsxtwxw/main.htm).
%%%%% This template is compatible with both Chinese and English.
%%%%% 正式编辑从line200开始
%%%%% Have fun with latex!


\documentclass[10pt,aspectratio=169]{beamer}

%%%%%%%%%%%%%%%%%%%%%%%%%%%%%%%%%%%%%%%%%%%%%%%%%%%%%%%%%%%%%%%
%%usepackage 所使用的宏包
%%%%%%%%%%%%%%%%%%%%%%%%%%%%%%%%%%%%%%%%%%%%%%%%%%%%%%%%%%%%%%%
\usepackage{t1enc}
\usepackage[utf8]{inputenc}
\usepackage{pgf,pgfarrows,pgfnodes,pgfautomata,pgfheaps,pgfshade}
\usepackage{graphicx}
\usepackage{amsmath,amsthm,amssymb,amsfonts}
\usepackage{color}
\usepackage{rotating}
\usepackage{multimedia}
\usepackage{natbib}
\usepackage{url}
\usepackage{ctex} 
\usepackage{lipsum,zhlipsum}
\usepackage{multicol}
\usepackage{tikz}
\usepackage{multirow}
\usepackage{booktabs} 
\usepackage{subfigure}
\usepackage{listings}
\usefonttheme[onlymath]{serif}
%%%%%%%%%%%%%%%%%%%%%%%%%%%%%%%%%%%%%%%%%%%%%%%%%%%%%%%%%%
%%define color 定义颜色
%%%%%%%%%%%%%%%%%%%%%%%%%%%%%%%%%%%%%%%%%%%%%%%%%%%%%%%%%%%%

\definecolor{SEU-green}{RGB}{77,124,44}
\definecolor{SEU-yellow}{RGB}{251,199,7}

%%%%%%%%%%%%%%%%%%%%%%%%%%%%%%%%%%%%%%%%%%%%%%%%%%%%%%%%%%%%%%%
%usetheme 所使用的主题
%%%%%%%%%%%%%%%%%%%%%%%%%%%%%%%%%%%%%%%%%%%%%%%%%%%%%%%%%%%%%%%
\usetheme[hideothersubsections]{Hannover}
%\usetheme[]{Pittsburgh}
%\usetheme[]{Warsaw}
%\usetheme[]{Copenhagen}
%\usetheme[]{Rochester}
%\usetheme[]{JuanLesPins}
%\usetheme[]{Berlin}
%\usetheme[]{Madrid}
%\useinnertheme{rounded}
%\usecolortheme{crane}
%\usecolortheme{dolphin}
%\usecolortheme{lily}
%\usecolortheme{beaver}

%\usefonttheme{structuresmallcapsserif}
\setbeamerfont{title}{family=\sf,size=\huge,series=\bfseries}
\usefonttheme{professionalfonts} % using non standard fonts for beamer
\usefonttheme{default} % default family is serif
\setbeamerfont{headline font} {family=\sf,size=\normalsize ,series=\bfseries}
\setbeamerfont{footline font} {family=\sf,size=\small, series=\bfseries}
\setbeamerfont{footnote} {family=\sf,size=\tiny,series=\bfseries}

%%%%%%%%%%%%%%%%%%%%%%%%%%%%%%%%%%%%%%%%%%%%%%%%%%%%%%%%%%%%%%%
%set font 设置英文字体 鸿蒙字体
%%%%%%%%%%%%%%%%%%%%%%%%%%%%%%%%%%%%%%%%%%%%%%%%%%%%%%%%%%%%%%%
\setmainfont[
Path=fonts/,
BoldFont=HarmonyOS_Sans_Bold.ttf,
ItalicFont=HarmonyOS_Sans_Regular_Italic.ttf,
BoldItalicFont=HarmonyOS_Sans_Bold_Italic.ttf,
]{HarmonyOS_Sans_Regular.ttf}
\setsansfont[
Path=fonts/,
BoldFont=HarmonyOS_Sans_Bold.ttf,
ItalicFont=HarmonyOS_Sans_Regular_Italic.ttf,
BoldItalicFont=HarmonyOS_Sans_Bold_Italic.ttf,
]{HarmonyOS_Sans_Regular.ttf}
\setmonofont[
Path=fonts/
]{HarmonyOS_Sans_Regular.ttf}

%%%%%%%%%%%%%%%%%%%%%%%%%%%%%%%%%%%%%%%%%%%%%%%%%%%%%%%%%%%%%%%
%set Chinese font 设置中文字体 鸿蒙字体
%%%%%%%%%%%%%%%%%%%%%%%%%%%%%%%%%%%%%%%%%%%%%%%%%%%%%%%%%%%%%%%
\setCJKmainfont[
Path=fonts/,
BoldFont=HarmonyOS_Sans_SC_Bold.ttf, 
ItalicFont=HarmonyOS_Sans_SC_Regular.ttf, 
BoldItalicFont=HarmonyOS_Sans_SC_Bold.ttf
]{HarmonyOS_Sans_SC_Regular.ttf}
\setCJKsansfont[
Path=fonts/,
BoldFont=HarmonyOS_Sans_SC_Bold.ttf, 
ItalicFont=HarmonyOS_Sans_SC_Regular.ttf, 
BoldItalicFont=HarmonyOS_Sans_SC_Bold.ttf
]{HarmonyOS_Sans_SC_Regular.ttf}
\setCJKmonofont[
Path=fonts/
]{HarmonyOS_Sans_SC_Regular.ttf}

%%%%%%%%%%%%%%%%%%%%%%%%%%%%%%%%%%%%%%%%%%%%%%%%%%%%%%%%%%%%%%%
%set beamer template
%%%%%%%%%%%%%%%%%%%%%%%%%%%%%%%%%%%%%%%%%%%%%%%%%%%%%%%%%%%%%%%
\setbeamertemplate{navigation symbols}{}
\setbeamertemplate{frametitle}[default][center]
\setbeamertemplate{itemize item}{\color{SEU-green}$\blacksquare$}
\setbeamertemplate{itemize subitem}{\color{SEU-yellow}$\blacktriangleright$}
\setbeamertemplate{itemize subsubitem}{\color{black}$\blacksquare$}
\setbeamertemplate{section}{\color{SEU-green}$\blacksquare$}

\setbeamercolor{title}{fg=SEU-green}
\setbeamercolor{author}{fg=black}
\setbeamercolor{frametitle}{fg=SEU-green}
\setbeamercolor{section in toc}{fg=SEU-green}
\setbeamercolor{title in sidebar}{fg=black}
\setbeamercolor{section in sidebar}{fg=red}
\setbeamercolor{section in sidebar shaded}{fg=gray}
\setbeamercolor{subsection in sidebar}{fg=red}
\setbeamercolor{subsection in sidebar shaded}{fg=gray}
\setbeamercolor{item projected}{bg=SEU-green}
\setbeamercolor{footnote}{fg=gray}

\beamertemplatenavigationsymbolsempty
\setbeamertemplate{caption}[numbered]

\setbeamertemplate{sidebar canvas left}[vertical shading][top=SEU-green!10,middle=SEU-green!10,bottom=SEU-green!10]



%default block
\setbeamercolor{block title}{fg=black,bg= SEU-green!50}
\setbeamercolor{block body}{fg=black,bg= SEU-green!5}

%example block
\setbeamercolor*{block body example}{fg=black,bg= SEU-yellow!5}
\setbeamercolor*{block title example}{fg=black,bg= SEU-yellow!50}

%alert block
\setbeamercolor{block title alerted}{bg=red!50,fg=black}
\setbeamercolor{block body alerted}{bg=red!5}

\AtBeginEnvironment{theorem}{%
	\setbeamercolor{block body}{bg=red!5}
	\setbeamercolor{block title}{bg=red!50,fg=black}
}


\setbeamertemplate{blocks}[default]%[rounded][shadow=ture]


\setbeamercolor{symbol color}{fg=SEU-green}
\setbeamercolor{structure}{fg=SEU-green}

%%%%%%%%%%%%%%%%%%%%%
\setbeamerfont{frametitle}{size=\LARGE,series=\bfseries}
\setbeamertemplate{frametitle}{
	\usebeamerfont{frametitle}
	\vspace{10pt}\insertframetitle	
}



\makeatletter
\addtobeamertemplate{sidebar left}{}{%

\hspace{0.25cm}	\includegraphics[width=1cm]{logo}%


\hspace{0.05cm}%

\begin{beamercolorbox}[leftskip=2 pt]{author in sidebar}{
		\href{mailto:XXXXXX@seu.edu.cn}{XXXXXX@seu.edu.cn}
}\end{beamercolorbox}
\vspace{0.05cm}%
	\begin{beamercolorbox}[leftskip=5 pt]{author in sidebar}{
	%\usebeamercolor[fg]{author in sidebar}
	%\usebeamerfont{author in sidebar}%
	%\insertshortauthor[width={1.5cm},center,respectlinebreaks]
	
	\insertdate 
}\end{beamercolorbox}

\hspace{0.05cm}

}


\makeatother


% 打开PDF后直接全屏
 \hypersetup{pdfpagemode=FullScreen}




%--------------------------------------------------------------%
% 正式编辑从line200开始，Formal editing starts at line200
%--------------------------------------------------------------%
\title[\fontsize{7pt}{0pt}\textbf{Short Title xxx研究}]{\vspace{2cm}\newline 硕士研究生论文答辩 \\\vspace{0.2cm} \textbf{XXXXXXXXXXXXXXXX研究}} 
\author[XXX \\ (English name)]{答辩人：XXX
%	\medskip
% \href{mailto:XXX@seu.edu.cn}{XXX@seu.edu.cn}
% 	\medskip
\\
校内导师：XXX
% 	\medskip
\\
校外导师：XXX}

\institute[东南大学]{东南大学XXX学院}
\date[ ]{\today}
\logo{\includegraphics[scale=0.25]{logo}}





\begin{document}
	{
%	\setbeamertemplate{headline}{\includegraphics[width=\paperwidth,height=\textheight,keepaspectratio=true]{head}}
\setbeamertemplate{background canvas}{%
	\includegraphics[width=\paperwidth,height=\paperheight,keepaspectratio=false]{head169}%
}
\setbeamertemplate{footline}{}
\setbeamertemplate{sidebar left}{ }
\setbeamertemplate{sidebar canvas left}{ }
\hspace*{-1.2cm}
%%%%%%%%%%%%%%%%%%
%\begin{frame}[plain,noframenumbering]
	\begin{frame}[noframenumbering]
	\centering 
	\titlepage
\end{frame}
}
% Background
\pgfdeclareimage[width=\paperwidth,height=\paperheight]{bg}{seu_background169.png}
\setbeamertemplate{background}{\pgfuseimage{bg}}





\begin{frame}[plain,noframenumbering]{目录}

	\tableofcontents
\end{frame}


\section[课题背景]{课题背景}

 \begin{frame}{课题背景}
		\begin{figure}  
		\includegraphics[height=0.5\paperheight]{example-image}  
		\caption{example-image\footnote{Donald K. Art of Computer Programming[M].  1973.
		}.}
\end{figure}



\end{frame}



 \begin{frame}{课题背景}
 双图页面
	\begin{multicols}{2}
		\begin{figure}  
			\includegraphics[height=0.5\paperheight]{example-image-a} 
			\caption{example-image-a.}
			\label{fig1}
		\end{figure}
		
		
		\begin{figure} 
			\includegraphics[height=0.5\paperheight]{example-image-b}  
			\caption{example-image-b.}
		\end{figure}
	\end{multicols}
	
\end{frame}



\section[研究内容]{开题报告研究内容}
\subsection{研究内容}
\begin{frame}{研究内容}
\qquad 针对xxxxxxxxxxxxxxxxx等问题，研究xxxxxxxxxxxxxxxxxxx技术。本课题按照“\alert{xxxx-xxxx-xxxx}”三个研究层级划分。3个研究内容分别为：


\begin{figure}
\includegraphics[width=0.6\paperwidth,keepaspectratio=T]{example-image-golden}
\end{figure}
\end{frame}



\subsection{技术路线}
\begin{frame}{技术路线}

\begin{figure}
	\includegraphics[width=0.4\paperwidth,keepaspectratio=T]{example-grid-100x100pt}
\end{figure}
\end{frame}


\begin{frame}{研究背景}  %开始一张幻灯片
	\begin{columns}[c]  %开始进入分栏环境，居中设置
		\column{0.5\paperwidth}  %第一栏（左栏）宽度为8cm
		\begin{itemize}  %栏内是一个列表
			\item{保险公司的经营是负债经营，其负债具有很大的不确定性。}
			\item{损失准备金是保险公司最重要的一项负债，其估计直接影响保险产品的定价以及最低资本的计算。}
			\item{1969年至2007年间破产的1023家保险公司中有38.1\%可归结为损失准备金提取不足(Coyne(2008))。}
			\item{损失准备金的估计历来是精算学研究的重要研究课题之一。}
		\end{itemize}
		\column{5cm}  %第二栏（右栏)宽度为3cm
		\begin{figure}  %栏内是一张图片
			\includegraphics[width=5cm,height = 5cm]{example-image-a}  
		\end{figure}
	\end{columns}  %分栏环境结束
\end{frame}




\begin{frame} %[<+->]   
%%如果需要逐条显示，删除上一行的“%”

  \frametitle{Basics}
  \begin{itemize}
  \item A frame starts with \textbackslash\texttt{begin\{frame\}}
    and end with \textbackslash\texttt{end\{frame\}}
  \item  Other environments inside a frame:
   \item For example \textbackslash\texttt{begin\{itemize\}}:
     \begin{itemize}
\item Whith a sign
\item[] Or not
\item \textbackslash\texttt{end\{itemize\}}
     \end{itemize}
   \end{itemize}
   \begin{enumerate}
\item Or \textbackslash\texttt{begin\{enumerate\}}:
\item which labels
\item the points with sequential numbers.
\item \textbackslash\texttt{end\{enumerate\}}
    \end{enumerate}
\end{frame}






\begin{frame}
	\frametitle{A Theorem on Infinite Sets}
	\begin{theorem}<1->
		There exists an infinite set.
	\end{theorem}
	\begin{proof}<2->
		This follows from the axiom of infinity.
	\end{proof}
	\begin{example}<3->[Natural Numbers]
		The set of natural numbers is infinite.
	\end{example}
\end{frame}







\section{研究进展}
\subsection{研究内容1}


\begin{frame}{研究内容1-子内容1}
\begin{multicols}{2}
	\lipsum[1]
		\begin{itemize}
		\item Lorem ipsum dolor sit amet
		\item Consectetuer adipiscing elit
		\item Ut purus elit, vestibulum ut
	\end{itemize}
\end{multicols}
	
\end{frame}


 \begin{frame}{研究内容1-子内容2}
	
			\zhlipsum[1][name = xiangyu]
	
\end{frame}



\begin{frame}{依次显示图片，本页可以自由控制图片在页面中所处的位置。}
			\begin{tikzpicture}[remember picture,overlay]   
		\node<1->[xshift=-3cm,yshift=-1cm] at (current page.center) {\includegraphics[height=4cm]{example-image-a}};  
	\node<2->[xshift=3cm,yshift=0cm] at (current page.center) {\includegraphics[height=4cm]{example-image-b}};
	\end{tikzpicture}
\end{frame}


\begin{frame}{左侧图片，右侧文字，另一种显示方法}
	\begin{picture}(100,50)
		\put(0,-50){\includegraphics[height=5cm]{example-image}}
		\put(200,80){\begin{minipage}[t]{0.5\linewidth}
				{Choose a point on the unit circle. Connect it to the origin with a line of length one, and denote the angle between that line and the horizontal coordinate axis by $\theta$.}
		\end{minipage}}
	\end{picture}
\end{frame}


\begin{frame}{在图片上悬浮加文字“$y=1.65298x+87.81814$”}
	\setlength{\unitlength}{1cm}
	
	\begin{figure}
		% Requires \usepackage{graphicx}
		\centering
		\includegraphics{example-image}%米氏常数拟合
		\begin{picture}(0,0)
			\put(-3,1){$y=1.65298x+87.81814$}
		\end{picture}
		\caption{过氧化氢酶$1/v\sim1/[\text{S}]$图}\label{figh2o2}
	\end{figure}

\end{frame}

\subsection{研究内容2}
\begin{frame}{Block}
	
块（block）共分为三种，分别是default、example和alert，颜色分别使用东大绿、东大黄、红色。

\begin{block}{默认块}
	  \begin{itemize}
  \item \textbackslash \texttt{setbeamercolor*{block body example}{fg=black, bg= SEU-green!5}}
\item \textbackslash \texttt{setbeamercolor*{block title example}{fg=black, bg=SEU-green!50}}
\item \textbackslash \texttt{setbeamertemplate{blocks}[rounded][shadow=true]}
  \end{itemize}
\end{block}

\begin{exampleblock}{例子}
	如果需要放入一些例子，可以用这个颜色的块。
\end{exampleblock}

\begin{alertblock}{提示块}
	这里展示的是需要额外注意的信息。
\end{alertblock}
\end{frame}
\subsection{研究内容3}
\begin{frame}
	\frametitle{List - open presentation mode to see how it works.}
	\begin{itemize}%[<+-| alert@+>]
		\pause
		\item To add one point at a time to a slide
		\pause
		\item you can use pause
		\begin{itemize}
			\pause
			\item part 1
			\pause
			\item part 2
		\end{itemize}
		\pause
		\item or alert:
		\pause
		\item \textbackslash\texttt{begin\{itemize\}[<+-| alert@+>]}
	\end{itemize}
\end{frame}
\subsection{研究内容4}

\begin{frame}{公式}
	\begin{equation}
		a^2+b^2=c^2
	\end{equation}
	
	\begin{equation}
		a^3+b^3=c^3
	\end{equation}
	\begin{equation}
		\mathcal{L}_{D}=\mathbb{E}_{x\sim P_{data}}[-\log D(x)]+\mathbb{E}_{z\sim P_{z}}[-\log (1-D(G(z)))]
		\label{equ-1}
	\end{equation}
\end{frame}






\begin{frame}
	\frametitle{\textbackslash\texttt{includegraphics} to insert figures.}
	\begin{columns}
		\column{0.5\textwidth}
		%\vspace{.5cm}
		%\begin{center}
		\includegraphics[scale=1.5]{logo}
		%\end{center}
		\column{0.6\textwidth}
		\scriptsize{\begin{itemize}
				\item \textbackslash\texttt{columns} can be used to devide the frame. 
				\item For two or more figures,
				\item or for a figure and some text.
		\end{itemize}}
	\end{columns}
\end{frame}





\section{总结}
\subsection{进度汇总}

\begin{frame}
  \frametitle{进度汇总}
% Table generated by Excel2LaTeX from sheet 'Sheet3'
\begin{table}[htbp]
	\centering
	\tiny
	\begin{tabular}{|c|l|}

		\toprule
	 \multicolumn{1}{|p{0.1\paperwidth}|}{\textbf{起讫日期}} & \textbf{工  作  内  容  和  要  求} \\
		\midrule
		\multicolumn{1}{|c|}{\multirow{4}[2]{*}{202x年xx月-202x年xx月}} & \textbf{工作内容：} \\
		& 收集xxxxxxxxxxxxxxxxxxxxxxxxxxxxxxxxxxxxxxxxxxxxxxxxxxxx。 \\
		& \textbf{成果目标：} \\
		& xxxxxxxxx。 \\
		\midrule
		\multicolumn{1}{|c|}{\multirow{4}[2]{*}{202x年xx月-202x年xx月}} & \textbf{工作内容：} \\
		& 建立xxxxx。 \\
		& \textbf{成果目标：} \\
		& 完成xxxx。 \\
		\midrule
		\multicolumn{1}{|c|}{\multirow{4}[2]{*}{202x年xx月-202x年xx月}} & \textbf{工作内容：} \\
		& 对比xxx。 \\
		& \textbf{成果目标：} \\
		& 完成xxx。 \\
		\midrule
		\multicolumn{1}{|c|}{\multirow{4}[2]{*}{202x年xx月-202x年xx月}} & \textbf{工作内容：} \\
		& 基于xxxxx。 \\
		& \textbf{成果目标：} \\
		& 完成xxxx；录用或发表论文1篇；申请专利1项。 \\
		\midrule
		\multicolumn{1}{|c|}{\multirow{4}[2]{*}{202x年xx月-202x年xx月}} & \textbf{工作内容：} \\
		& 研究xxxxx。 \\
		& \textbf{成果目标：} \\
		& 提出xxxx，申请专利1项。 \\
		\midrule
		\multicolumn{1}{|c|}{\multirow{4}[2]{*}{202x年xx月-202x年xx月}} & \textbf{工作内容：} \\
		& 研究xxxxxxxxxxxxxxxxxxxxxxxxxxxxxxxxxxxxxxxxxxxxxxxxxxxxxxxxxxxxxxxxxxxxxxxxxxxxxxxxxxxxxxxxxxxx。 \\
		& \textbf{成果目标：} \\
		& 提出xxx，录用或发表论文1篇。 \\
		\midrule
		\multicolumn{1}{|c|}{\multirow{4}[2]{*}{202x年xx月-202x年xx月}} & \textbf{工作内容：} \\
		& xxxxx。 \\
		& \textbf{成果目标：} \\
     	& 梳理课题研究成果，整理并提交硕士论文，准备答辩。 \\
		\bottomrule
	
	\end{tabular}%
	\label{tab:addlabel}%
\end{table}%



\begin{tikzpicture}[overlay,remember picture]
	\node(1)[xshift=2.5cm,yshift=0.8cm] at (current page.center) {\includegraphics[scale=0.5]{fig/全部完成.pdf}};
\end{tikzpicture}

\end{frame}

\subsection{已取得的阶段成果}
\begin{frame}
	\frametitle{已取得的阶段成果}
	(1)论文
	\begin{itemize}
		\item [1]	\textbf{xxx}, xxxx*. xxxxxxxxxxxxxxxxxxx[J]. xxxxxxxxxx, 2021, x: xxxxx. 
		(第一作者, SCI检索)
		
		\item [2]	\textbf{xxx}, xxxx*. xxxxxxxxxxxxxxxxxxx[J]. xxxxxxxxxx, 2021, x: xxxxx. 
		(第一作者, SCI检索)
	\end{itemize}
(2) 专利
	\begin{itemize}
	\item [1]	xxxx (第二发明人，授权)
	\item[2]	xxxx (第二发明人，受理)
	\item [3]	xxxx (第二发明人，受理)
	\item [4]	xxxx (第一发明人，在投)
\end{itemize}
	
	
	
\end{frame}

	{
	\setbeamertemplate{background canvas}{%
		\includegraphics[width=\paperwidth,height=\paperheight,keepaspectratio=false]{head169}
	}
	\setbeamertemplate{footline}{}
	\logo{\includegraphics[scale=0.25]{logo}\hspace{\dimexpr\paperwidth-1.8cm}\vspace{-0cm}}
	\setbeamertemplate{background}{}
		\setbeamertemplate{sidebar left}{ }
	\setbeamertemplate{sidebar canvas left}{ }
	\hspace*{-1.2cm}
	\frame{
		\vspace{3cm}
		
		\begin{beamercolorbox}{title}
			\usebeamerfont{title}
		\centering \Huge \bfseries
		欢迎批评指正！
	\end{beamercolorbox}
\vspace{2cm}
\begin{beamercolorbox}{author}
	\centering
	\insertauthor \\
	\vspace{2ex}
	\insertinstitute
\end{beamercolorbox}
}
}


 \end{document}






%%% Local Variables: 
%%% mode: latex
%%% TeX-master: t
%%% End: 
